\begin{thebibliography}{99}
\raggedright
\addcontentsline{toc}{section}{References}

\bibitem{repo}
Vladimir Reshetnikov, \emph{ProveIt: Topology/UnknotRecognition}, source
commit \texttt{fdeb5b1a20b84b93cb150e0232031837bb7d30cb},
9 October 2026 UTC.
\href{https://github.com/VladimirReshetnikov/ProveIt/tree/fdeb5b1a20b84b93cb150e0232031837bb7d30cb/Topology/UnknotRecognition}{Pinned repository tree}.
In particular, the maintained \texttt{fast} implementation, tests,
research drivers, and accompanying \texttt{synthesis} sections.

\bibitem{aht}
Ian Agol, Joel Hass, and William Thurston,
\emph{The computational complexity of knot genus and spanning area},
Transactions of the American Mathematical Society \textbf{358}(9)
(2006), 3821--3850.
\href{https://arxiv.org/abs/math/0205057v2}{arXiv:math/0205057v2}.
The orbit-counting and weighted results are Theorems 12 and 16 in the
referenced preprint version.

\bibitem{lackenby-talk}
Marc Lackenby, \emph{Unknot recognition in quasi-polynomial time},
Oxford talk handout, March 2021.
\href{https://people.maths.ox.ac.uk/lackenby/quasipolynomial-talk-oxford-compressed.pdf}{Author's handout}.
The announced $2^{O((\log n)^3)}$ bound appears on slide 7;
the hierarchy-depth outline appears on slide 25.

\bibitem{lackenby-2026}
Marc Lackenby, \emph{Incompressible surfaces, hierarchies and unknot
recognition}, preprint, 25 July 2026.
\href{https://arxiv.org/abs/2607.23350v1}{arXiv:2607.23350v1}.
See Section 9 for the complexity discussion and Proposition 9.1.

\bibitem{lackenby-curves}
Marc Lackenby, \emph{Some fast algorithms for curves in surfaces},
Discrete \& Computational Geometry \textbf{76} (2026), 539--588.
\href{https://doi.org/10.1007/s00454-026-00845-7}{doi:10.1007/s00454-026-00845-7}.

\bibitem{schleimer}
Saul Schleimer, \emph{Polynomial-time word problems},
Commentarii Mathematici Helvetici \textbf{83}(4) (2008), 741--765.
\href{https://doi.org/10.4171/CMH/142}{doi:10.4171/CMH/142}.
\href{https://arxiv.org/abs/math/0608563}{arXiv:math/0608563}.

\bibitem{kozen}
D. Kozen, U. V. Vazirani, and V. V. Vazirani,
\emph{NC algorithms for comparability graphs, interval graphs, and
testing for unique perfect matching}, in Foundations of Software Technology and
Theoretical Computer Science, Lecture Notes in Computer Science
\textbf{206} (1985), 496--503.
\href{https://doi.org/10.1007/3-540-16042-6_28}{doi:10.1007/3-540-16042-6\_28}.
Expanded as \emph{NC Algorithms for Comparability Graphs, Interval Graphs,
and Unique Perfect Matchings}, Cornell TR 86-799 (1986).
\href{https://www.cs.cornell.edu/~kozen/Papers/comp.pdf}{Authors' manuscript}.
The source matching/acyclicity characterization is classical.

\bibitem{golumbic}
M. C. Golumbic, T. Hirst, and M. Lewenstein,
\emph{Uniquely restricted matchings}, Algorithmica \textbf{31}(2)
(2001), 139--154.
\href{https://doi.org/10.1007/s00453-001-0004-z}{doi:10.1007/s00453-001-0004-z}.

\bibitem{dowling}
William F. Dowling and Jean H. Gallier,
\emph{Linear-time algorithms for testing the satisfiability of
propositional Horn formulae}, Journal of Logic Programming
\textbf{1}(3) (1984), 267--284.
\href{https://doi.org/10.1016/0743-1066(84)90014-1}{doi:10.1016/0743-1066(84)90014-1}.
\href{https://www.seas.upenn.edu/~cis5110/Dowling-Gallier-Horn-sat.pdf}{Authors' article copy}.

\bibitem{wirtinger}
Ryan Blair, Alexandra Kjuchukova, Roman Velazquez, and Paul Villanueva,
\emph{Wirtinger systems of generators of knot groups},
Communications in Analysis and Geometry \textbf{28}(2) (2020), 243--262.
\href{https://doi.org/10.4310/CAG.2020.v28.n2.a2}{doi:10.4310/CAG.2020.v28.n2.a2}.
\href{https://arxiv.org/abs/1705.03108}{arXiv:1705.03108}.

\bibitem{incompatibility}
Ryan Blair, Alexandra Kjuchukova, and Makoto Ozawa,
\emph{The incompatibility of crossing number and bridge number for knot
diagrams}, Discrete Mathematics \textbf{342}(7) (2019), 1966--1978.
\href{https://doi.org/10.1016/j.disc.2019.03.013}{doi:10.1016/j.disc.2019.03.013}.
\href{https://arxiv.org/abs/1710.11327}{arXiv:1710.11327}.

\bibitem{erickson-nayyeri}
Jeff Erickson and Amir Nayyeri,
\emph{Tracing compressed curves in triangulated surfaces},
Discrete \& Computational Geometry \textbf{49}(4) (2013), 823--863.
\href{https://jeffe.cs.illinois.edu/pubs/pdf/tracing.pdf}{Authors' manuscript}.
Section 3.3 discusses the compressed tracing history.

\bibitem{papakyriakopoulos}
C. D. Papakyriakopoulos,
\emph{On Dehn's lemma and the asphericity of knots},
Proceedings of the National Academy of Sciences of the United States
of America \textbf{43}(1) (1957), 169--172.
\href{https://doi.org/10.1073/pnas.43.1.169}{doi:10.1073/pnas.43.1.169}.
Theorem 3(i) states the free-cyclic knot-group characterization of the
unknot. The longer proof appears under the same title in Annals of
Mathematics \textbf{66}(1) (1957), 1--26.

\bibitem{reapr}
Jason Cantarella, Henrik Schumacher, and Clayton Shonkwiler,
\emph{Hard unknots are often easy from a different perspective},
preprint, version 3, October 2026.
\href{https://arxiv.org/abs/2607.28772v3}{arXiv:2607.28772v3}.
The associated test sets are available at
\href{https://doi.org/10.5061/dryad.vx0k6dk6v}{doi:10.5061/dryad.vx0k6dk6v}.

\end{thebibliography}
